\chapter{Bewegung}
Bevor die Umsetzung der Erfassung geklärt werden kann, muss definiert werden, was unter Bewegung zu verstehen ist, um die optimale Stelle am Körper für die Datenerfassung zu finden.

\section{Definition der Bewegung in Bezug auf das System}
In dieser Arbeit werden drei verschiedene Bewegungsarten analysiert: Stehen, Gehen und Rennen. Jede dieser Bewegungsarten zeichnet sich durch spezifische biomechanische Merkmale aus. Beim Gehen bewegt sich der Körper analog zu einem umgekehrten Pendel – die kinetische Energie der ersten Standphase wird in potenzielle Gravitationsenergie umgewandelt und anschließend teilweise zurückgewonnen, wenn der Körper nach vorne und unten fällt \cite{cappellini2006motor}. Im Gegensatz dazu ähnelt das Rennen dem Hüpfen auf einem Pogo-Stick, bei dem kinetische und potenzielle Gravitationsenergie vorübergehend als elastische Dehnungsenergie in Muskeln, Sehnen und Bändern gespeichert und in der zweiten Hälfte der Standphase nahezu vollständig zurückgewonnen wird \cite{cappellini2006motor}. Diese Unterschiede ist auf die Abbildung \ref{fig:walking_running} zu erkennen.

Ein weiterer bedeutender Unterschied zwischen diesen Bewegungen liegt in den Phasen des Bodenkontakts: Während das Gehen durch eine doppelte Stützphase gekennzeichnet ist, in der beide Füße kurzzeitig den Boden berühren, weist das Rennen eine „Flugphase“ auf, in der sich keine der Gliedmaßen am Boden befindet \cite{cappellini2006motor}. Diese Unterschiede sind für die spätere Klassifizierung der Bewegungen essenziell, da sie als zentrale Merkmale zur Unterscheidung genutzt werden können.

\begin{figure}[h]
	\centering
	\includegraphics[width=0.5\linewidth]{walkingrunning}
	\caption{Walking and Running  \cite{cappellini2006motor}}
	\label{fig:walking_running}
\end{figure}

\section{Welche Körperteile werden genutzt?}
Um die Unterschiede zwischen Stehen, Gehen und Rennen zu analysieren, ist es wichtig zu verstehen, welche Körperteile beteiligt sind, um die optimale Platzierung des Sensors zu bestimmen. Zu Beginn der Beschleunigung sind die Quadrizeps, also die Oberschenkelmuskeln, am stärksten beteiligt. Der Musculus soleus und der Musculus gastrocnemius, beides Skelettmuskeln der Unterschenkelmuskulatur, spielen ebenfalls eine wichtige Rolle. Obwohl die Arme nur wenig zur direkten Vorwärtsbewegung oder zum Abstützen beitragen, sind sie für den Ausgleich der vertikalen Drehbewegung der unteren Extremitäten von entscheidender Bedeutung. \cite{hamner2010muscle} Man könnte daher annehmen, dass die Platzierung des Sensors am Oberschenkel und am Handgelenk die beste Wahl ist.

Um diese Hypothese zu überprüfen und die optimale Sensorplatzierung zu bestimmen, bietet die Studie von Özdemir \cite{ozdemir2016analysis} wertvolle Einblicke. Die Studie analysiert die Sensorplatzierung von Bewegungserkennungssystemen und liefert wertvolle Informationen für die Entwicklung tragbarer Geräte zur Überwachung menschlicher Bewegungen. Es wurden verschiedene Sensorplatzierungen am Kopf, an der Brust, an der Taille, am rechten Handgelenk, am rechten Oberschenkel und am rechten Unterschenkel untersucht. Die Genauigkeit wurde dann mit maschinellen Lernverfahren untersucht, wie wir es später auch tun werden.

Das Ergebnis zeigt, dass die Taille die am besten geeignete Stelle für die Platzierung des Sensors am Körper ist, während der Sensor am Handgelenk die schlechtesten Ergebnisse lieferte, obwohl dies die bevorzugte Position für viele tragbare Sensoren ist. Für die Erfassung von Bewegungen wie Stehen, Gehen oder Rennen ist daher die Platzierung des Sensors an der Taille die beste Wahl. Die Taille befindet sich in der Nähe des Körperschwerpunkts und kann daher die Bewegungen des gesamten Körpers besser erfassen. Andere Körperstellen, wie z.B. das Handgelenk, sind weniger gut geeignet, wie die Studie gezeigt hat. Aus diesem Grund haben wir uns für die Messung an der Taille entschieden, wie in Abbildung \ref{fig:sensor} dargestellt.

\begin{figure}[h]
	\centering
	\includegraphics[width=0.5\linewidth]{sensor}
	\caption{Mögliche Sensorplatzierungen (eigene Abbildung)}
	\label{fig:sensor}
\end{figure}
